% TJ93090C
% Pneumothorax
% 14.0
% 2013-11-30

.. _m-pneumothorax:


Taktik \Ca{Zeitkritisch! Vitale Bedrohung!}
Entlastungspunktion, zügiger Transport
in eine geeignete Einrichtung}

\begin{itemize}
  -   |TxMassVitMKBes|
  \begin{itemize}
    -   **Lagerung**: Oberkörper hoch. Wenn stabile
    Seitenlage notwendig ist, dann Lagerung auf die
    *verletzte* Seite, damit die unverletzte
    Lunge  nicht unnötig belastet wird.
    -   \Z{**Notarzt**: Ein Pneumothorax kann mittels
    Drainage bzw. Punktion entlastet werden (ärztliche Maßnahme).}
  \end{itemize}
  -   **Verband** bei offener Thoraxverletzung: **Nicht
  luftdicht verbinden**, es kann sonst ein
  Spannungspneumothorax entstehen! Steril abdecken.
  -   **Transportentscheidung**: Schockraum
\end{itemize}